% ===== Preâmbulo compartilhado (Tectonic / XeTeX) — documentos da FTL094 =====
\documentclass[11pt,a4paper]{article}
\usepackage{fontspec}
\usepackage[brazil]{babel}
\usepackage{geometry}
\geometry{a4paper,margin=2.5cm}
\usepackage{amsmath,amssymb}
\usepackage{graphicx}
\usepackage{booktabs}
\usepackage{array}
\usepackage{longtable}
\usepackage{tabularx}
\usepackage{float}
\usepackage{enumitem}
\usepackage{xcolor}
\usepackage{tikz}
\usetikzlibrary{arrows.meta,positioning,automata,shapes.geometric,fit,backgrounds,calc}
\usepackage{pgfplots}
\pgfplotsset{compat=1.18}
\usepackage{listings}
\usepackage{caption}
\usepackage{fancyhdr}
\usepackage[hidelinks]{hyperref}

% ----- cores -----
\definecolor{ufamblue}{RGB}{0,51,102}
\definecolor{codegreen}{rgb}{0,0.5,0}
\definecolor{codeblue}{rgb}{0,0,0.7}
\definecolor{codered}{rgb}{0.6,0,0}
\definecolor{lgray}{rgb}{0.95,0.95,0.95}

% ----- listagens C -----
\lstdefinestyle{c}{language=C,basicstyle=\ttfamily\footnotesize,
  keywordstyle=\color{codeblue}\bfseries,commentstyle=\color{codegreen}\itshape,
  stringstyle=\color{codered},numbers=left,numberstyle=\tiny\color{gray},
  frame=single,breaklines=true,tabsize=2,showstringspaces=false,columns=fullflexible}

% ----- constantes do projeto -----
\newcommand{\projnome}{Sistema de Navegação Autônoma para o Robô Móvel Khepera~IV}
\newcommand{\disciplina}{FTL094 --- Engenharia de Software de Tempo Real}
\newcommand{\periodo}{Período Letivo 2026/1}
\newcommand{\versaodoc}{1.0}

% ----- cabeçalho/rodapé -----
\pagestyle{fancy}\fancyhf{}
\rhead{\small \tipodoc}
\lhead{\small Khepera~IV --- FTL094}
\cfoot{\thepage}
\renewcommand{\headrulewidth}{0.4pt}

% ----- coluna X centralizada p/ tabularx -----
\newcolumntype{C}{>{\centering\arraybackslash}X}

% ----- capa padronizada -----
\newcommand{\fazcapa}{%
\begin{titlepage}\centering
{\large\bfseries UNIVERSIDADE FEDERAL DO AMAZONAS (UFAM)}\\[3pt]
{\normalsize Centro de P\&D em Tecnologia Eletrônica e da Informação --- CETELI}\\[2pt]
{\normalsize Programa de Pós-Graduação em Engenharia Elétrica --- PPGEE}\\[26pt]
{\large \disciplina}\\[2pt]
{\normalsize \periodo}\\[64pt]
\rule{\textwidth}{1.2pt}\\[12pt]
{\Huge\bfseries\color{ufamblue} \tipodoc}\\[14pt]
{\Large \projnome}\\[10pt]
\rule{\textwidth}{1.2pt}\\[48pt]
{\large\bfseries Equipe}\\[8pt]
{\normalsize Gabriel Henrique de Souza Nazaré}\\[3pt]
{\normalsize Luan Alexandre Ramos Barreto}\\[3pt]
{\normalsize Matheus Rocha Canto}\\[3pt]
{\normalsize Matheus Serrão Uchôa}\\[3pt]
{\normalsize Ricardo Mendonça Braz}\\[3pt]
{\normalsize Yan Fernandes da Silva}\\[40pt]
{\normalsize Docente responsável: Prof. Dr.~[nome do docente]}\\
\vfill
{\small Documento de tipo \emph{\tipodoc} --- Versão \versaodoc}\\[2pt]
{\small Manaus --- AM, \datadoc}
\end{titlepage}}

% ----- tabela de controle de versões -----
\newcommand{\controleversoes}{%
\section*{Controle de versões}
\begin{tabularx}{\textwidth}{@{}l l l X@{}}
\toprule
\textbf{Versão} & \textbf{Data} & \textbf{Autor} & \textbf{Descrição} \\
\midrule
1.0 & \datadoc & Equipe FTL094 & Emissão inicial do documento. \\
\bottomrule
\end{tabularx}
\bigskip}

\newcommand{\tipodoc}{Arquitetura e Projeto do Sistema}
\newcommand{\datadoc}{30 de junho de 2026}

\begin{document}
\fazcapa
\controleversoes
\tableofcontents
\newpage

% =====================================================================
\section{Introdução}
% =====================================================================
\subsection{Propósito e escopo}
Este documento descreve a \textbf{arquitetura} e o \textbf{projeto detalhado} do \emph{\projnome}, derivando-os dos requisitos especificados no documento de Análise de Requisitos. Cobre o estilo arquitetural, a decomposição em módulos, o comportamento dinâmico (máquina de estados e laço de controle), o projeto algorítmico e as decisões de projeto com sua justificativa.

\subsection{Requisitos arquiteturalmente significativos}
Influenciam a arquitetura, sobretudo: RNF-01 (segurança/sem colisão), RNF-04 (portabilidade sim$\to$real), RT-01--RT-03 (tempo real) e RF-06 (escape de mínimo local). Estes conduzem a um estilo \emph{reativo em camadas} com isolamento da localização.

% =====================================================================
\section{Visão arquitetural}
% =====================================================================
\subsection{Estilo adotado}
Adota-se uma \textbf{arquitetura híbrida em \emph{pipeline} (Sentir--Decidir--Atuar) com arbitragem por máquina de estados}, combinando uma camada \emph{deliberativa} mínima (atração ao alvo) e uma camada \emph{reativa} (desvio por proximidade). O estilo \emph{pipe-and-filter} organiza o fluxo de dados percepção$\to$decisão$\to$atuação; a arbitragem por FSM implementa a ideia de subsunção (o desvio subsome a atração quando há obstáculo próximo).

\subsection{Diagrama de componentes}
\begin{figure}[H]\centering
\begin{tikzpicture}[
  comp/.style={draw,rounded corners,minimum height=1.1cm,minimum width=2.6cm,align=center,font=\small,fill=ufamblue!8},
  io/.style={draw,minimum height=1.1cm,minimum width=2.0cm,align=center,font=\small,fill=gray!12},
  >={Stealth},node distance=0.9cm and 1.0cm]
  \node[io] (sens) {Sensores\\IR/GPS/IMU};
  \node[comp,right=of sens] (perc) {Percepção e\\Estimação de\\estado};
  \node[comp,right=of perc] (fsm) {Arbitragem\\(FSM)};
  \node[comp,right=of fsm] (ctrl) {Controle e\\Cinemática};
  \node[io,right=of ctrl] (mot) {Atuadores\\(motores)};
  \node[comp,below=1.1cm of fsm] (log) {Telemetria\\(logging)};
  \draw[->] (sens)--(perc); \draw[->] (perc)--(fsm); \draw[->] (fsm)--(ctrl); \draw[->] (ctrl)--(mot);
  \draw[->] (perc.south) |- (log.west);
  \draw[->] (fsm.south) -- (log.north);
\end{tikzpicture}
\caption{Componentes do software e fluxo de dados do laço de controle.}
\label{fig:comp}
\end{figure}

% =====================================================================
\section{Decomposição estrutural}
% =====================================================================
\begin{table}[H]\centering
\caption{Módulos, responsabilidades e requisitos atendidos.}
\renewcommand{\arraystretch}{1.2}
\begin{tabularx}{\textwidth}{@{}l X l@{}}
\toprule
\textbf{Módulo} & \textbf{Responsabilidade} & \textbf{Requisitos} \\
\midrule
Percepção & Ler e condicionar sensores; estimar pose $(x,y,\theta)$ e ``massas'' de obstáculo. & RF-04, RF-02 \\
Arbitragem (FSM) & Decidir o estado (Navegar/Girar/Contornar) e o comando de alto nível. & RF-02, RF-06 \\
Controle & Lei de rumo (P) + termo reativo; converter $(v,\omega)$ em rodas. & RF-01, RF-05 \\
Atuação & Aplicar velocidades aos motores; parada segura. & RF-03, RF-05 \\
Telemetria & Registrar periodicamente o estado interno em arquivo. & RF-07, RNF-06 \\
Configuração & Centralizar alvo e parâmetros (constantes). & RF-08, RNF-05 \\
\bottomrule
\end{tabularx}
\end{table}

\subsection{Comportamento dinâmico: máquina de estados}
\begin{figure}[H]\centering
\begin{tikzpicture}[->,>=Stealth,node distance=3.4cm,on grid,
  every state/.style={draw,rounded corners,minimum width=2.2cm,align=center,font=\small}]
  \node[state] (N) {\textsc{Navegar}\\\scriptsize vai ao alvo};
  \node[state] (G) [right=of N] {\textsc{Girar}\\\scriptsize gira até limpar};
  \node[state] (C) [right=of G] {\textsc{Contornar}\\\scriptsize passa a caixa};
  \path
    (N) edge[bend left] node[above,font=\scriptsize]{frente $>$ \textsc{trig}} (G)
    (G) edge[bend left] node[above,font=\scriptsize]{frente $<$ \textsc{clear}} (C)
    (C) edge[bend left] node[below,font=\scriptsize]{nova parede} (G)
    (C) edge[bend left=55] node[below,font=\scriptsize]{percorreu $d_p$} (N);
\end{tikzpicture}
\caption{Máquina de estados de arbitração (atende RF-02 e RF-06).}
\label{fig:fsm}
\end{figure}

\subsection{Laço de controle (visão de atividade)}
\begin{figure}[H]\centering
\begin{tikzpicture}[node distance=0.7cm,>={Stealth},
  act/.style={draw,rounded corners,minimum width=3.4cm,minimum height=0.8cm,align=center,font=\small},
  dec/.style={draw,diamond,aspect=2,align=center,font=\scriptsize,inner sep=1pt}]
  \node[act] (t) {Aguarda tique ($T=16$~ms)};
  \node[act,below=of t] (p) {Lê sensores / estima pose};
  \node[dec,below=of p] (d) {chegou?};
  \node[act,below=of d] (f) {Avalia FSM e calcula $(v,\omega)$};
  \node[act,below=of f] (a) {Aplica velocidades + telemetria};
  \node[act,right=2.2cm of d] (stop) {Para motores (fim)};
  \draw[->] (t)--(p); \draw[->] (p)--(d);
  \draw[->] (d)-- node[right,font=\scriptsize]{não} (f);
  \draw[->] (d)-- node[above,font=\scriptsize]{sim} (stop);
  \draw[->] (f)--(a);
  \draw[->] (a.west) -- ++(-2.0,0) |- (t.west);
\end{tikzpicture}
\caption{Atividade do laço de controle periódico (tarefa de tempo real).}
\label{fig:loop}
\end{figure}

% =====================================================================
\section{Projeto detalhado}
% =====================================================================
\subsection{Estado interno e parâmetros}
O estado de execução compreende a pose $(x,y,\theta)$, o estado da FSM, a direção de escape e o ponto de início do contorno. Os parâmetros (alvo, ganhos $K_\theta,K_a$, limiares de IR, $d_p$, período $T$) são constantes centralizadas (RNF-05).

\subsection{Leis de controle}
A distância e o erro de rumo ao alvo:
\begin{align}
  d &= \sqrt{(x_g-x)^2+(y_g-y)^2}, &
  e_\theta &= \operatorname{wrap}\!\big(\operatorname{atan2}(y_g-y,x_g-x)-\theta\big).
\end{align}
No estado \textsc{Navegar}, o comando de giro soma atração (P de rumo) e repulsão reativa:
\begin{equation}
  \omega = K_\theta\,e_\theta + K_a\,(b_R-b_L),
\end{equation}
com $b_L,b_R$ as massas de obstáculo à esquerda/direita estimadas dos IR acima da linha de base $s_0$. A conversão para rodas (tração diferencial):
\begin{equation}
  \dot\phi_L = v-\omega, \qquad \dot\phi_R = v+\omega.
\end{equation}
Nos estados \textsc{Girar}/\textsc{Contornar}, o robô fixa a direção de escape e a mantém (compromisso de contorno) até liberar a frente e ultrapassar o obstáculo por $d_p$.

\subsection{Excerto do projeto da FSM}
\begin{lstlisting}[style=c,caption={Núcleo da máquina de estados.}]
switch (state) {
  case NAVEGAR:
    steer = K_HEADING*err + K_AVOID*(block_r - block_l);
    forward = CRUISE * slowdown(f, err);
    if (front > IR_TRIG) { escape_dir = (block_l<=block_r)?+1:-1; state = GIRAR; }
    break;
  case GIRAR:
    forward = 0; steer = escape_dir*SPIN;
    if (f<IR_CLEAR && fl<IR_CLEAR && fr<IR_CLEAR) { mark(&pass); state = CONTORNAR; }
    break;
  case CONTORNAR:
    forward = 0.7*CRUISE; steer = K_AVOID*(block_r - block_l);
    if (front > IR_TRIG) state = GIRAR;
    else if (dist_from(pass) > PASS_DIST) state = NAVEGAR;
    break;
}
\end{lstlisting}

\subsection{Aspectos de tempo real}
O laço é uma \textbf{tarefa periódica} de período $T=16$~ms e \emph{deadline} $D=T$ (RT-01, RT-02). O projeto evita alocação dinâmica e operações de custo variável no caminho crítico, mantendo o WCET muito inferior a $T$ (RT-03). Em uma futura execução sobre RTOS, o conjunto monotarefa é trivialmente escalonável; havendo tarefas concorrentes (p.ex.\ visão), recomenda-se análise por taxa monotônica (RMS).

% =====================================================================
\section{Decisões de projeto e justificativa}
% =====================================================================
\begin{table}[H]\centering
\caption{Registro de decisões de arquitetura (ADR).}
\renewcommand{\arraystretch}{1.25}
\begin{tabularx}{\textwidth}{@{}l X X@{}}
\toprule
\textbf{Id} & \textbf{Decisão} & \textbf{Justificativa / alternativas} \\
\midrule
AD-01 & Linguagem \textbf{C}. & Tempo real e toolchain embarcada no Webots; alternativa Python descartada (sem interpretador instalado, menor previsibilidade). \\
AD-02 & Arquitetura \textbf{híbrida} reativa+deliberativa. & IR de curto alcance e ausência de mapa tornam planejamento global inviável; puramente reativo é frágil. \\
AD-03 & \textbf{FSM com contorno comprometido} (Bug/\emph{wall-following}). & Resolve o mínimo local do campo potencial puro (RF-06). \\
AD-04 & \textbf{GPS+IMU} em simulação. & Isola navegação de localização (RNF-04); plano de migração para odometria no robô real. \\
AD-05 & \textbf{Calibração empírica} dos sensores. & A leitura real diverge da folha de dados (reflexão do piso); medir é mandatório. \\
\bottomrule
\end{tabularx}
\end{table}

\section{Rastreabilidade arquitetura $\leftrightarrow$ requisitos}
Os componentes da Figura~\ref{fig:comp} cobrem os requisitos conforme a Tabela~2: Percepção (RF-04, RF-02), FSM (RF-02, RF-06), Controle/Atuação (RF-01, RF-03, RF-05), Telemetria (RF-07, RNF-06); o laço periódico (Figura~\ref{fig:loop}) atende RT-01--RT-03; a estratégia de localização isolada atende RNF-04.

\end{document}
